\section{历年题：题型方法与双层解答}

\begin{methodbox}[title={题型：证明连续像紧致}]
看到“\(X\) 紧致、\(f\) 连续、证明 \(f(X)\) 紧致”，应立即按紧致性的定义处理任意开覆盖。
构造只有一个自然方向：把像空间的覆盖取原像，借连续性得到定义域的开覆盖，再用紧致性
选出有限个。若题目只要求这一结论，不需要引入序列紧致等更强且可能不适用的工具。
\end{methodbox}

\problemstart
\subsection{示例真题（\problemid{SAMPLE-XXXX-X-Q01}）}

\paragraph{题目}
设 \(X\) 为紧致拓扑空间，\(f:X\to Y\) 连续。证明 \(f(X)\) 紧致。

\begin{examanswer}
任取 \(f(X)\) 的开覆盖 \(\{V_i\}_{i\in I}\)。对每个 \(i\)，取 \(Y\) 中开集
\(W_i\)，使 \(V_i=f(X)\cap W_i\)。由 \(f\) 连续，
\(\{f^{-1}(W_i)\}_{i\in I}\) 是 \(X\) 的开覆盖。

由 \(X\) 紧致，存在 \(i_1,\dots,i_m\in I\)，使
\[
  X=\bigcup_{k=1}^{m}f^{-1}(W_{i_k}).
\]
于是对每个 \(y=f(x)\in f(X)\)，某个 \(k\) 满足
\(x\in f^{-1}(W_{i_k})\)，从而 \(y\in V_{i_k}\)。故
\(f(X)=\bigcup_{k=1}^{m}V_{i_k}\)，所以 \(f(X)\) 紧致。
\end{examanswer}

\begin{idea}
题目要证明紧致，定义要求我们从“任意开覆盖”出发。因此第一步不是寻找特殊点或序列，
而是把 \(f(X)\) 的任意开覆盖摆在桌面上。现在唯一未使用的条件是连续性；连续性最擅长
把开集通过原像拉回去，所以自然地考虑 \(f^{-1}(W_i)\)。

拉回后的覆盖位于紧致空间 \(X\) 上，因而能选出有限子覆盖。最后要检查有限个原像覆盖
\(X\) 是否真的推出对应的有限个 \(V_i\) 覆盖 \(f(X)\)：任取 \(y\in f(X)\)，选择
原像 \(x\)，再沿有限覆盖找到包含 \(x\) 的原像开集即可。这一“覆盖拉回—有限化—推回”
链条就是应当记住的可重建结构。定义和完整证明见
\knowledgelink{sec:compactness}{紧致性与连续性}。
\end{idea}
